\chapter{Études de cas II}\label{ch:fm:case-studies-ii}

Le 1\up{er} août 2012, un teneur de marché, Knight Capital, perd plus de 460~millions de dollars en quarante-cinq minutes. Cinq jours plus tard, la société a vendu pour 400~millions de dollars d'actions privilégiées convertibles à un groupe d'investisseurs,
convertibles en environ 73~\% de ses actions~; onze mois plus tard, elle fusionne avec GETCO au sein d'une nouvelle
société holding. Les contrôles ont échoué le matin~; les décisions qui ont fixé le coût de la défaillance pour les propriétaires de la société sont intervenues dans les cinq jours suivants, et
bien avant pour certaines. Ce chapitre examine quatre autres défaillances tirées des données publiques~: un incident logiciel, les prime brokers d'un family office, une
bourse qui était aussi une contrepartie, et la décision d'une bourse face à une tension sur le marché.

\section{2012~: un incident logiciel et un sauvetage}

La décision de la Securities and Exchange Commission datée d'octobre 2013 retrace la matinée~: un technicien n'a pas copié une nouvelle version du routeur d'ordres de
la société sur l'un de ses huit serveurs, où une ancienne fonction subsistait et restait appelable. Pendant 45~minutes, ce routeur a envoyé des millions d'ordres et obtenu plus de 4~millions d'exécutions sur 154~actions,
pour plus de 397~millions d'actions. Le livre~11 a examiné cette décision sous l'angle de ses contrôles~: aucune limite automatique sur l'exposition globale de la société, et une surveillance qui
n'a pas déclenché d'alerte. La société a perdu plus de 460~millions de dollars sur les chiffres retenus par cette décision~; son rapport trimestriel évalue la perte réalisée avant impôt à environ 440~millions de dollars,
pour des capitaux propres totaux de 1\,497~millions de dollars au 30~juin.

\begin{dated}{2026-09}{Les données publiques des cas du chapitre}\label{dat:fm:case-studies-ii:record}
\textbf{Knight Capital} (décision SEC 34-70694~; formulaires 10-Q, 8-K du 6 août 2012 et 8-K12G3 de KCG du 1\up{er} juillet 2013)~: plus de 460~millions de dollars perdus
le 1\up{er} août 2012~; 400~millions de dollars d'actions privilégiées convertibles vendues le 6 août, convertibles en environ 73~\% des actions ordinaires, au titre de l'exception de la bourse
en matière de viabilité financière, dérogeant à l'approbation des actionnaires~; fusion avec GETCO au sein de KCG Holdings le 1\up{er} juillet 2013, chaque action échangée contre 3.75~dollars en espèces
ou un tiers d'action KCG. \textbf{Archegos} (rapport du comité spécial de Credit Suisse du 29 juillet 2021)~: Credit Suisse a perdu environ 5.5~milliards de dollars~; le Conseil
de la Réserve fédérale a infligé à UBS Group une amende de 268.5~millions de dollars pour la gestion du risque de contrepartie de Credit Suisse le 24 juillet 2023, soit environ 387~millions de dollars avec la sanction de la
Prudential Regulation Authority~; le fondateur a été condamné en juillet 2024 à 18~ans de prison. \textbf{Nickel du LME}~: la Divisional Court
a jugé que la bourse avait agi légalement en annulant les transactions du 8 mars 2022, et la Court of Appeal a rejeté le recours le 7 octobre 2024~; la FCA
a infligé à la bourse une amende de 9\,245\,900~livres le 19 mars 2025 pour ne pas avoir maintenu un marché ordonné.
\end{dated}

Le 6 août, la société a vendu 400\,000 actions privilégiées convertibles pour 400~millions de dollars à un groupe comprenant Jefferies, Blackstone et GETCO,
convertibles en environ 266.7~millions d'actions ordinaires à 1.50~dollar chacune, pour 97.9~millions d'actions en circulation. Une émission de cette ampleur exige normalement un
vote des actionnaires~; le comité des finances et d'audit du conseil s'est appuyé à la place sur l'exception de la bourse pour un retard qui mettrait gravement en péril la
viabilité financière de la société, et la bourse ne s'y est pas opposée. Onze mois plus tard, la société a fusionné avec GETCO au sein d'une nouvelle société holding, et
chaque action est devenue 3.75~dollars en espèces ou un tiers d'action de la nouvelle société.

\begin{table}[ht]
\centering
\begin{tabular}{>{\raggedright\arraybackslash}p{5.6cm}>{\centering\arraybackslash}p{2.6cm}>{\centering\arraybackslash}p{2.6cm}>{\centering\arraybackslash}p{2.6cm}}
\toprule
 & 30 juin 2012 & après la perte & après le sauvetage \\
\midrule
Capitaux propres comptables (millions de dollars) & 1\,497 & 1\,057 & 1\,457 \\
Actions (millions) & 97.9 & 97.9 & 364.6 \\
Part des anciens actionnaires & 100~\% & 100~\% & 26.9~\% \\
Capitaux propres comptables des anciens actionnaires (millions de dollars) & 1\,497 & 1\,057 & 391 \\
Valeur comptable par action (dollars) & 15.29 & 10.80 & 4.00 \\
\bottomrule
\end{tabular}
\caption{Le sauvetage de 2012 en valeurs comptables, tiré du rapport trimestriel et du contrat de cession, avant impôt~: la perte a absorbé 29.4~\% des capitaux propres, le
sauvetage à 1.50~dollar par action ayant emporté la plus grande part de ce qui restait aux anciens actionnaires. Données~: \texttt{fm\_cases2.knight}.}\label{tab:fm:case-studies-ii:rescue}
\end{table}

Le tableau est l'arithmétique, au niveau de la société, d'un sauvetage conduit sous la pression du temps. La perte seule a coûté aux actionnaires 29.4~\% de leurs capitaux propres comptables~; le sauvetage,
fixé à 1.50~dollar par action alors que la valeur comptable par action était encore de 10.80~dollars, leur a coûté 666~millions de dollars supplémentaires de capitaux propres comptables, transférés aux nouveaux
investisseurs, qui ont payé 400~millions de dollars pour 1\,066~millions de dollars. Le prix a été fixé par les jours qui restaient à la société~: ses clients et ses contreparties,
selon le rapport trimestriel, avaient réduit leurs flux d'ordres et perdu confiance, et une société sans capital ne peut pas continuer à faire du market making. Ce qu'elle n'avait pas décidé
à l'avance (la limite sur l'exposition globale, le capital détenu face à une matinée de ce genre, une facilité permanente qui pouvait être tirée en une journée) a fixé les
conditions.

\begin{omfigure}
\begin{tikzpicture}
\begin{axis}[width=9.5cm,height=5.2cm,ybar stacked,bar width=26pt,ymin=0,ymax=1700,ylabel={\small capitaux propres comptables (millions de dollars)},xtick={0,1,2},
  xticklabels={30 juin 2012,après la perte,après le sauvetage},xmin=-0.6,xmax=2.6,tick label style={font=\footnotesize},grid=major,
  grid style={gray!20},table/col sep=comma,legend style={font=\scriptsize,at={(0.5,-0.2)},anchor=north,
  /tikz/every even column/.append style={column sep=8pt}},legend columns=2,legend cell align=left]
\addplot[fill=omDef!60,draw=omDef] table[x=pos,y=old]{figdata/desk/29-case-studies-ii/knight.csv};
\addplot[fill=omThm!45,draw=omThm] table[x=pos,y=new]{figdata/desk/29-case-studies-ii/knight.csv};
\coordinate (k0) at (axis cs:0,1497);
\coordinate (k1) at (axis cs:1,1057);
\coordinate (k2) at (axis cs:2,196);
\coordinate (k3) at (axis cs:2,924);
\legend{anciens actionnaires,investisseurs du sauvetage}
\end{axis}
\node[font=\scriptsize,above] at (k0) {1\,497};
\node[font=\scriptsize,above] at (k1) {1\,057};
\node[font=\scriptsize] at (k2) {391};
\node[font=\scriptsize] at (k3) {1\,066};
\end{tikzpicture}\par
\omcaption{Qui détenait les capitaux propres comptables du teneur de marché en 2012~: la perte, puis le sauvetage à 1.50~dollar par action. Données~: \texttt{fm\_cases2.knight}.}\label{fig:fm:case-studies-ii:knight}
\end{omfigure}

\section{2021~: un family office et ses prime brokers}

Le comité spécial du conseil de Credit Suisse a publié son rapport sur les pertes de la banque liées à Archegos, un family office, en juillet 2021. Le soir du
25 mars 2021, après que la valeur de ses plus grosses positions s'est effondrée depuis le 22 mars, Archegos a indiqué à ses prime brokers qu'il avait 120~milliards de dollars d'exposition brute,
dont 70~milliards de dollars à l'achat et 50~milliards à la vente, pour 9 à 10~milliards de dollars de \omterm{def:fm:the-proprietary-market-making-firm:ownfunds}{fonds propres}, et leur a demandé d'accepter un moratoire pendant qu'il liquidait
ses positions. Ils ont refusé. L'exposition brute propre de Credit Suisse s'élevait à environ 27~milliards de dollars le 23 mars et à 17~milliards le 26 mars~; elle avait appelé
2.8~milliards de dollars de marge le 25 mars~; elle a perdu environ 5.5~milliards de dollars.

\begin{definition}[Appels de marge dynamiques]\label{def:fm:case-studies-ii:dynamic}
Les \emph{appels de marge dynamiques}\index{appels de marge dynamiques} fixent la marge d'un client à partir du risque courant de son portefeuille (sa volatilité, sa concentration,
sa liquidité et son biais directionnel), recalculée à mesure que les positions et les prix évoluent, plutôt que comme un pourcentage fixe du notionnel de chaque transaction à l'ouverture.
\end{definition}

Le rapport décrit les deux régimes au sein d'une même banque. Le book de prime brokerage était soumis à des \omterm{def:fm:case-studies-ii:dynamic}{appels de marge dynamiques}~; les swaps étaient margés de façon statique, la marge initiale
étant fixée en dollars à l'ouverture, de sorte qu'à mesure que les positions prenaient de la valeur, la marge rapportée à leur montant s'érodait. La marge moyenne sur les swaps s'élevait
à environ 9.4~\% le 23 mars. Une position concentrée que son propre détenteur estimait à deux semaines ou un mois pour la liquider portait la marge
d'une position diversifiée.

\begin{definition}[Course à la sortie]\label{def:fm:case-studies-ii:race}
Une \emph{course à la sortie}\index{course à la sortie} est la liquidation séquentielle, par plusieurs créanciers, de collatéraux ou de couvertures portant sur les mêmes positions concentrées
après le défaut d'un client commun~: chacun vend dans des prix déjà déplacés par ceux qui ont vendu avant lui, de sorte que la perte de chaque créancier dépend de son rang dans l'ordre
autant que de sa marge.
\end{definition}

Le rapport retrace la course. Le 26 mars, Goldman a organisé des ventes de blocs portant sur une partie des positions, et Credit Suisse a participé à trois d'entre elles sans savoir
combien d'actions étaient offertes ni lesquelles avaient été vendues en premier~; elle a vendu un peu plus de 3~milliards de dollars ce jour-là, en maintenant ses ventes algorithmiques entre 2 et 3~\% du
volume quotidien moyen. Pendant le week-end, plusieurs banques, dont Deutsche Bank, Morgan Stanley et Goldman, n'ont pas montré d'intérêt pour une
liquidation gérée~; Credit Suisse, UBS et Nomura en ont accepté une et ont vendu des blocs en avril.

\section{Tutoriel~: le sauvetage, la course, la ruée, la tension}\label{tut:fm:case-studies-ii:tut}

\begin{tutorial}
\textbf{Objectif.} Étendre le recueil de cas aux quatre cas~: la dilution du sauvetage de 2012~; une \omterm{def:fm:case-studies-ii:race}{course à la sortie} entre sept prime brokers calibrée sur le rapport~;
le déficit et la ruée sur la bourse de 2022~; et l'appel de marge d'une vente à découvert sur le nickel face à ses liquidités. \textbf{État final~:} le
\cref{tab:fm:case-studies-ii:rescue} et les \cref{fig:fm:case-studies-ii:knight,fig:fm:case-studies-ii:race,fig:fm:case-studies-ii:ftx,fig:fm:case-studies-ii:nickel}.
\begin{tutsteps}
\item \textbf{Le sauvetage.} \texttt{firm.casebook.rescue} avec les capitaux propres, la perte, 400~millions de dollars pour 266.7 millions d'actions, et 97.9 millions d'anciennes actions.
\item \textbf{La course.} Sept courtiers (Credit Suisse et les six que nomme le rapport) détiennent chacun un septième de la position, pour 17~milliards de dollars chacun comme Credit
Suisse le 26 mars, soit environ 120~milliards de dollars d'exposition brute au total. Ils vendent à tour de rôle~; le prix baisse linéairement avec la fraction cumulée vendue (impact
permanent du livre~10).
\item \textbf{La calibration.} Le coefficient d'impact est fixé de sorte qu'un courtier disposant de l'exposition de Credit Suisse et de sa marge de 9.4~\%, vendant en sixième, perde
5.5~milliards de dollars~: \texttt{race\_impact} donne 0.531, soit une baisse de 53~\% une fois que les sept ont vendu. Le rang est une hypothèse~: le rapport place Credit
Suisse parmi les courtiers qui ont vendu en avril.
\item \textbf{La ruée.} \texttt{ftx\_balances} et \texttt{ftx\_flows} lisent la transcription que fait le livre~3 du tableau des déficits des débiteurs.
\item \textbf{La tension.} Une vente à découvert à titre illustratif de 10\,000 tonnes de nickel, adossée à 300~millions de dollars de liquidités, margée à partir du cours de clôture du 7 mars.
\end{tutsteps}
\end{tutorial}

\omcode{code/firm/casebook/firm_casebook.py}{101}{116}{La course à la sortie~: les courtiers vendent à tour de rôle sous un impact permanent linéaire, chacun perdant la baisse sur son
segment, nette de sa marge~; et le coefficient d'impact qui reproduit une perte donnée.}\label{lst:fm:case-studies-ii:race}

\begin{omfigure}
\begin{tikzpicture}
\begin{axis}[width=10.5cm,height=5.6cm,xlabel={\small rang dans la course à la sortie},ylabel={\small perte du courtier (milliards de \$)},xmin=0.7,xmax=7.3,
  ymin=-0.3,ymax=7.8,xtick={1,2,3,4,5,6,7},tick label style={font=\footnotesize},grid=major,grid style={gray!20},table/col sep=comma,
  legend cell align=left,legend style={font=\scriptsize,at={(0.03,0.97)},anchor=north west}]
\addplot[omThm,thick,mark=triangle*] table[x=place,y=m075]{figdata/desk/29-case-studies-ii/race.csv};
\addplot[omDef,very thick,mark=*] table[x=place,y=m094]{figdata/desk/29-case-studies-ii/race.csv};
\addplot[gray,thick,dashed,mark=square*] table[x=place,y=m200]{figdata/desk/29-case-studies-ii/race.csv};
\draw[black,thin] (axis cs:6,5.5) circle[radius=3pt];
\node[font=\scriptsize,anchor=south east] at (axis cs:5.9,5.7) {calibration};
\legend{marge 7.5~\%,marge 9.4~\%,marge 20~\%}
\end{axis}
\end{tikzpicture}\par
\omcaption{Une \omterm{def:fm:case-studies-ii:race}{course à la sortie} stylisée entre sept prime brokers, chacun avec 17~milliards de dollars d'exposition sur les mêmes positions, calibrée de sorte que la sixième,
avec une marge de 9.4~\%, perde 5.5~milliards de dollars. Le premier vendeur ne perd rien~; le dernier perd 6.79~milliards de dollars à 9.4~\% et 4.99~milliards de dollars à 20~\%. Données~:
\texttt{fm\_cases2.race}.}\label{fig:fm:case-studies-ii:race}
\end{omfigure}

Dans le modèle, le premier courtier à vendre ne perd rien à aucune des trois marges~: il vend à une baisse moyenne de 3.8~\%, à l'intérieur de sa marge. Le dernier
perd 6.79~milliards de dollars à 9.4~\%, et 4.99~milliards de dollars même à 20~\%. La marge qui aurait laissé chaque courtier indemne est la baisse moyenne sur son
segment~: 3.8~\% pour le premier, 26.6~\% pour le milieu, 49.3~\% pour le dernier. Aucun pourcentage statique ne couvre la course, car la perte dépend
de l'ordre, que nul courtier ne contrôle à l'avance.

\begin{proposition}[La course redistribue une baisse fixe]\label{prop:fm:case-studies-ii:zero}
Avec un impact permanent linéaire $1-\kappa f$ et des courtiers détenant les fractions $s_1,\dots,s_K$ de la position, dont la somme vaut un, la baisse totale de valeur
réalisée par l'ensemble des courtiers vaut $\kappa/2$ de la position, quel que soit l'ordre~; l'ordre décide seulement de qui la supporte. Le courtier au rang $k$ supporte
$s_k\kappa\bigl(F_{k-1}+s_k/2\bigr)$, où $F_{k-1}$ est la fraction vendue avant lui.
\end{proposition}

\begin{proof}
Le courtier $k$ vend entre les fractions cumulées $F_{k-1}$ et $F_k=F_{k-1}+s_k$ au prix moyen $1-\kappa(F_{k-1}+s_k/2)$, si bien que sa baisse vaut
$s_k\kappa(F_{k-1}+s_k/2)=\kappa\int_{F_{k-1}}^{F_k}f\,df$. En sommant sur $k$, on obtient $\kappa\int_0^1f\,df=\kappa/2$.
\end{proof}

La proposition explique pourquoi un statu quo est rationnel pour les créanciers pris ensemble et irrationnel pour chacun pris isolément~: une liquidation gérée laisse la baisse totale
inchangée (ou plus petite, si des ventes plus lentes réduisent l'impact), alors que vendre en premier reporte sur les autres sa propre part de celle-ci. Les \omterm{def:fm:case-studies-ii:dynamic}{appels de marge dynamiques}
sont la défense dont dispose chaque courtier~: une marge qui croît avec la taille de la position rapportée au volume du marché et avec l'ignorance du courtier
quant à ce que détient le client ailleurs, fixée avant le début de la course.

\section{2022~: une bourse qui était aussi une contrepartie}

Le livre~3 retrace l'effondrement de FTX à partir des dossiers des débiteurs. Pour une société de trading, la leçon porte sur l'exposition. Une société qui détient des liquidités et des jetons sur une bourse
est créancière de cette bourse~; si la bourse prête les actifs de ses clients à une entité affiliée, la société est aussi, sans le savoir, créancière de cette
entité affiliée. Le premier rapport intermédiaire des débiteurs établit que cette entité, Alameda, disposait d'une limite d'emprunt fixée à 65~milliards de dollars et de paramètres lui permettant de retirer en deçà de cette limite et de laisser un solde négatif~; l'analyse du déficit
établit qu'à la date de la requête, pour les jetons les plus liquides, la bourse détenait des actifs
valant 6.6~\% de ce qu'elle devait à ses clients, et 22.6~\% pour l'ensemble des jetons assortis de créances.

\begin{omfigure}
\begin{tikzpicture}
\begin{axis}[width=10.5cm,height=5.4cm,xlabel={\small novembre 2022},ylabel={\small flux net cumulé (milliards de \$)},xmin=1,xmax=11,ymin=-7.8,ymax=4,
  xtick={1,2,3,4,5,6,7,8,9,10,11},tick label style={font=\footnotesize},grid=major,grid style={gray!20},table/col sep=comma,legend cell align=left,
  legend style={font=\scriptsize,at={(0.03,0.97)},anchor=north west}]
\addplot[omDef,very thick,mark=*,mark size=1.5pt] table[x=day,y=customers]{figdata/desk/29-case-studies-ii/ftx.csv};
\addplot[omThm,thick,dashed,mark=square*,mark size=1.5pt] table[x=day,y=related]{figdata/desk/29-case-studies-ii/ftx.csv};
\draw[gray,dotted] (axis cs:8,-7.8) -- (axis cs:8,4) node[pos=0.3,right,font=\scriptsize] {retraits suspendus};
\draw[black,thin] (axis cs:1,0) -- (axis cs:11,0);
\legend{clients,Alameda et parties liées}
\end{axis}
\end{tikzpicture}\par
\omcaption{La ruée sur la bourse de 2022~: les retraits nets cumulés des clients ont atteint 7.0~milliards de dollars à la date de la requête, dont 3.3~milliards de dollars rien que le 7 novembre
tandis que l'entité affiliée et les parties liées ont déposé un montant net de 3.2~milliards de dollars. Données~: \texttt{fm\_cases2.ftx}, d'après la transcription que fait le livre~3 du
tableau des débiteurs.}\label{fig:fm:case-studies-ii:ftx}
\end{omfigure}

La ruée a été rapide~: les clients ont retiré un montant net de 3.3~milliards de dollars le 7 novembre, et le 8 novembre la bourse a suspendu les retraits. Une société qui a découvert le
déficit le 8 novembre n'avait plus aucune issue~; les décisions qui comptaient étaient les limites fixées plus tôt sur le montant détenu sur une seule plateforme de négociation, la
rapidité avec laquelle les soldes excédentaires étaient rapatriés, et la nature de la preuve de ségrégation qu'elle exigeait (les \omterm{def:fm:crisis-management:trapped}{actifs bloqués} du chapitre~27).

\section{2022~: la tension sur le nickel et la décision de la bourse}

Le 7 mars 2022, le contrat à trois mois sur le nickel du London Metal Exchange a clôturé à 48\,078~dollars la tonne, en hausse depuis une ouverture à peine inférieure à 30\,000~dollars. En
les premières heures du 8 mars, il est monté à 101\,365~dollars à 06~h~08 et s'est maintenu au-dessus de 80\,000~dollars à partir de 07~h~00. La décision de la Court of Appeal consigne ce à quoi la
bourse faisait face~: si les transactions du matin étaient maintenues, les membres auraient dû déposer environ 19.75~milliards de dollars de marge d'ici 09~h~00, et au moins cinq
devaient faire défaut selon les prévisions. Elle a suspendu les transactions à 08~h~15 et annulé chaque transaction conclue depuis minuit. De grandes positions courtes avaient été constituées par plusieurs
intervenants, parmi lesquels Tsingshan sur le marché de gré à gré. L'avis de la FCA ajoute que les appels de marge avaient établi des records de 3.5~milliards de dollars le 4
mars et 5.1~milliards de dollars le matin du 7 mars, et que les bandes de prix avaient été désactivées de 04~h~49 à 08~h~15 sur décision de collaborateurs juniors.

\begin{omfigure}
\begin{tikzpicture}
\begin{axis}[width=10.5cm,height=5.4cm,xlabel={\small cours du nickel (milliers de \$ la tonne)},ylabel={\small appel de marge de variation (millions de \$)},
  xmin=46,xmax=104,ymin=0,ymax=600,tick label style={font=\footnotesize},grid=major,grid style={gray!20},table/col sep=comma]
\addplot[omDef,very thick] table[x=price,y=call]{figdata/desk/29-case-studies-ii/nickel.csv};
\draw[omThm,thick,dashed] (axis cs:46,300) -- (axis cs:104,300) node[pos=0.02,above right,font=\scriptsize] {liquidités disponibles};
\draw[gray,dotted] (axis cs:78.078,0) -- (axis cs:78.078,600) node[pos=0.93,left,font=\scriptsize] {liquidités épuisées};
\draw[gray,dotted] (axis cs:101.365,0) -- (axis cs:101.365,600) node[pos=0.93,left,font=\scriptsize] {pic};
\end{axis}
\end{tikzpicture}\par
\omcaption{Une vente à découvert à titre illustratif de 10\,000 tonnes de nickel, margée à partir du cours de clôture du 7 mars~: son appel atteint 319~millions de dollars à 80\,000~dollars et 533
millions de dollars au pic, et ses 300~millions de dollars de liquidités s'épuisent à 78\,078~dollars. Données~: \texttt{fm\_cases2.nickel}.}\label{fig:fm:case-studies-ii:nickel}
\end{omfigure}

Pour la vente à découvert, l'arithmétique est celle du chapitre~14~: un appel qui croît linéairement avec le cours, face à des liquidités qui ne croissent pas.
Pour la bourse, il s'agissait d'un choix entre un appel de marge qui aurait fait défaut à plusieurs membres et une annulation qui aurait retiré les transactions à ceux qui les avaient conclues.
La Divisional Court a jugé l'annulation légale et a rejeté la demande similaire de Jane Street~; la Court of Appeal a rejeté l'appel d'Elliott, en estimant que
la seule autre solution, percevoir la marge au cours de clôture précédent, aurait placé la chambre de compensation en manquement à ses propres obligations. Deux enseignements
au niveau de la société en découlent pour une société de trading~: une bourse peut annuler des transactions dans des conditions extrêmes, de sorte qu'un profit réalisé sur un marché désordonné n'est pas encore un profit~;
et le plan de marge de la société doit résister à des appels plusieurs fois supérieurs à tout ce qui a été vu auparavant, puisque les appels records sont survenus lors de séances consécutives.

\section{Ce que les quatre cas ont en commun}

\begin{method}[D'un cas à un contrôle]\label{met:fm:case-studies-ii:control}
\begin{enumerate}
\item Énoncer la décision que le dossier montre avoir été prise, ou non, avant l'événement~: une limite, une méthode de marge, une exposition sur une plateforme de négociation, un plan de financement.
\item Dimensionner l'événement à partir des chiffres publics~: la perte rapportée aux \omterm{def:fm:the-proprietary-market-making-firm:ownfunds}{fonds propres}, les pertes de la \omterm{def:fm:case-studies-ii:race}{course à la sortie} par place, le déficit, l'appel face aux liquidités.
\item Rédiger le contrôle qui aurait modifié cette dimension~: une limite d'exposition agrégée et un plan de capital pour y faire face~; des \omterm{def:fm:case-studies-ii:dynamic}{appels de marge dynamiques} assortis de
majorations pour concentration~; des limites et des balayages sur les soldes des plateformes~; un plan de marge pour les appels records.
\item Tester le contrôle sur le portefeuille de la société à l'aide du modèle du recueil de cas, et consigner le résultat dans le registre des risques (chapitre~12).
\end{enumerate}
\end{method}

Les quatre défaillances diffèrent par leur cause et partagent une même structure. Dans chacune, les conditions d'une contrepartie avaient été fixées avant l'événement, sur la base d'hypothèses que l'événement a mises en défaut~: un
sauvetage chiffré en quelques jours, une marge fixée à l'ouverture, un solde laissé sur une plateforme, un appel de marge calibré sur des mouvements ordinaires. Et dans chacun des cas, les sociétés qui s'en sont le mieux tirées
avaient arrêté leur réponse à l'avance~: les courtiers qui ont vendu les premiers, la bourse dotée du pouvoir d'annuler, les sociétés dont les expositions étaient déjà
réduites. L'exercice du chapitre~27 permet à une société de découvrir laquelle de ces décisions elle a prises.

\section{Construction~: le recueil de cas, partie 2}\label{bld:fm:case-studies-ii:casebook}

\begin{build}
\bfield{Rôle} Les chiffres publics des quatre cas sous forme de tableaux liés à leurs lignes du grand livre~; un calculateur de dilution lors d'un sauvetage~; la \omterm{def:fm:case-studies-ii:race}{course à la sortie} avec une calibration
sur une perte annoncée~; le déficit et la ruée de la bourse à partir des données du livre~3~; l'appel de marge d'une vente à découvert face à ses liquidités.
\bfield{Interface} \texttt{firm.casebook}~: \texttt{KNIGHT}, \texttt{ARCHEGOS}, \texttt{NICKEL}, \texttt{rescue}, \texttt{exit\_race},
\texttt{race\_impact}, \texttt{ftx\_balances}, \texttt{ftx\_flows}, \texttt{margin\_call}, \texttt{break\_price}.
\bfield{Règles} Les chiffres proviennent du grand livre ou de données versionnées accompagnées d'une ligne de licence~; les données illustratives sont désignées comme telles~; la baisse totale de la course est
indépendante de l'ordre (la \cref{prop:fm:case-studies-ii:zero}).
\bfield{Tests d'acceptation} \texttt{code/firm/casebook/tests/}~: la propriété de somme nulle de la course et sa calibration, le sauvetage sur des chiffres calculés à la main, le
calcul de l'appel de marge et du prix en cas d'échec, et la couverture du tableau des déficits.
\bfield{Pour aller plus loin} Impact en racine carrée dans la course, avec la vitesse de vente propre à chaque courtier~; une liquidation gérée sous la forme d'un calendrier commun~; la vente à découvert sur le nickel
avec une marge initiale relevée par la chambre de compensation.
\end{build}

\begin{omsources}
\item US Securities and Exchange Commission, In the Matter of Knight Capital Americas LLC, Release 34-70694, 16 octobre 2013.
\item Knight Capital Group, Form 10-Q for the quarter ended 30 June 2012; Form 8-K, 6 août 2012; KCG Holdings, Form 8-K12G3, 1 juillet 2013.
\item Credit Suisse Group Special Committee of the Board of Directors, \emph{Report on Archegos Capital Management}, 29 juillet 2021.
\item Board of Governors of the Federal Reserve System, communiqué de presse, 24 juillet 2023~; US Attorney's Office, SDNY, communiqué de presse, 19 décembre 2024.
\item FTX Debtors, \emph{Preliminary Analysis of Shortfalls}, 2 mars 2023~; J.~J. Ray III, First Interim Report, 9 avril 2023.
\item R (Elliott Associates) v London Metal Exchange [2024] EWCA Civ 1168; Financial Conduct Authority, Final Notice to the London Metal Exchange,
19 mars 2025.
\end{omsources}

\section{Exercices}

\begin{exercise}[$\star$]\label{exo:fm:case-studies-ii:1}
À partir du \cref{tab:fm:case-studies-ii:rescue}, calculer la fraction de la société de 2012 que détenaient encore ses anciens actionnaires après le sauvetage, et
le prix par action payé par les investisseurs du sauvetage~?
\end{exercise}

\begin{exercise}[$\star$]\label{exo:fm:case-studies-ii:2}
Définir les \omterm{def:fm:case-studies-ii:dynamic}{appels de marge dynamiques} et expliquer pourquoi une marge statique sur une position concentrée dont le cours monte s'érode.
\end{exercise}

\begin{exercise}[$\star$]\label{exo:fm:case-studies-ii:3}
Une vente à découvert de 10\,000 tonnes est margée à partir de 48\,078~dollars~: calculer son appel à 80\,000~dollars, et le prix auquel les 300~millions de dollars de liquidités s'épuisent~?
\end{exercise}

\begin{exercise}[$\star\star$]\label{exo:fm:case-studies-ii:4}
Démontrer la \cref{prop:fm:case-studies-ii:zero}. Qu'implique-t-elle pour un accord de standstill entre créanciers~?
\end{exercise}

\begin{exercise}[$\star\star$]\label{exo:fm:case-studies-ii:5}
Dans la course calibrée, déterminer la marge qui laisse indemnes le premier, le sixième et le dernier courtier, et celle qui couvre la position
si un seul courtier la détenait en totalité~?
\end{exercise}

\begin{exercise}[$\star\star$]\label{exo:fm:case-studies-ii:6}
Déterminer quelles limites auraient protégé une société de trading disposant de soldes sur la bourse de 2022, et quels registres lui auraient signalé que ces soldes étaient menacés~?
\end{exercise}

\begin{exercise}[$\star\star\star$]\label{exo:fm:case-studies-ii:7}
\emph{Programmation.} Relancer la course avec une marge de 20~\%. Quelles places ne perdent rien, et quelle perte subit la dernière~?
\end{exercise}

\begin{exercise}[$\star\star\star$]\label{exo:fm:case-studies-ii:8}
\emph{Trouver la faille.} «~Notre marge sur le client est de 20~\%, soit le double de celle du secteur, donc nous ne pouvons pas y perdre.~»
\end{exercise}

\section{Problème~: Le dernier à vendre}

\begin{problem}[{Problème du week-end --- le dernier à vendre}]\label{pb:fm:case-studies-ii:1}
Le responsable de l'activité de prime brokerage doit expliquer au conseil d'administration ce que signifient les quatre cas pour sa politique de marge et pour ses propres expositions.

\medskip\noindent\textbf{Partie I --- Les notions.}
\begin{enumerate}
\item Définir les \omterm{def:fm:case-studies-ii:dynamic}{appels de marge dynamiques}.
\item Définir une \omterm{def:fm:case-studies-ii:race}{course à la sortie}.
\item Énoncer la \cref{prop:fm:case-studies-ii:zero}.
\item Pourquoi un standstill est-il rationnel pour les créanciers pris ensemble et non pris isolément~?
\end{enumerate}

\medskip\noindent\textbf{Partie II --- Le dossier.}
\begin{enumerate}[resume]
\item Résumer le dossier public de l'incident de 2012 et le sauvetage qui en a découlé.
\item Indiquer ce que le family office a dit à ses prime brokers le 25 mars 2021, et ce qu'ils ont fait.
\item Résumer le déficit et la ruée de la bourse de 2022.
\item Indiquer ce qu'a décidé la bourse du marché du nickel, et ce que les tribunaux et la FCA ont établi.
\end{enumerate}

\medskip\noindent\textbf{Partie III --- Les reconstitutions.}
\begin{enumerate}[resume]
\item Donner le tableau de dilution du sauvetage.
\item Décrire le modèle de \omterm{def:fm:case-studies-ii:race}{course à la sortie} et sa calibration.
\item Donner les pertes par place pour une marge de 9.4~\%.
\item Donner les pertes par place pour une marge de 20~\%, ainsi que la marge dont chaque place a besoin.
\item Indiquer ce que le modèle laisse de côté.
\item Donner les appels de la vente à découvert illustrative sur le nickel et le prix en cas d'échec.
\end{enumerate}

\medskip\noindent\textbf{Partie IV --- La politique.}
\begin{enumerate}[resume]
\item Quelle méthode de marge adopteriez-vous pour les clients concentrés~?
\item Comment sauriez-vous ce qu'un client détient chez d'autres courtiers~?
\item Quelles limites fixeriez-vous sur les soldes propres de la société auprès des plateformes~?
\item Que modifieriez-vous dans le plan de capital de la société après le cas de 2012~?
\item Énoncer le \emph{résultat nommé}~: la perte de chaque prime broker en fonction de son rang dans la \omterm{def:fm:case-studies-ii:race}{course à la sortie} et de sa marge, calibrée sur les chiffres publics,
et la dilution du sauvetage de 2012.
\item En deux phrases, rédiger la recommandation.
\end{enumerate}
\end{problem}

\section{Questions d'entretien}

\begin{interviewq}[$\star$ \iqroles{risque}]\label{iq:fm:case-studies-ii:1}
Quelle est la différence entre appels de marge statiques et \omterm{def:fm:case-studies-ii:dynamic}{appels de marge dynamiques}~?
\end{interviewq}

\begin{interviewq}[$\star$ \iqroles{développeur, risque}]\label{iq:fm:case-studies-ii:2}
Une nouvelle version commence à envoyer des ordres que personne n'avait prévus. Quels contrôles auraient dû l'arrêter, et quelle décision détermine ce qu'elle coûte aux propriétaires~?
\end{interviewq}

\begin{interviewq}[$\star\star$ \iqroles{risque, trader}]\label{iq:fm:case-studies-ii:3}
Un client fait défaut sur des swaps que vous et cinq autres banques détenez sur les mêmes actions. Les autres proposent un standstill. Que faites-vous~?
\end{interviewq}

\begin{interviewq}[$\star\star$ \iqroles{risque}]\label{iq:fm:case-studies-ii:4}
Quelle part de la trésorerie de la société conserveriez-vous sur une seule plateforme de crypto-actifs, et pourquoi~?
\end{interviewq}

\begin{interviewq}[$\star\star$ \iqroles{trader}]\label{iq:fm:case-studies-ii:5}
Vous avez réalisé un gain important sur une transaction dans un marché que la bourse a ensuite déclaré désordonné. À quoi faut-il s'attendre~?
\end{interviewq}

\begin{interviewq}[$\star\star\star$ \iqroles{chercheur, risque}]\label{iq:fm:case-studies-ii:6}
Modélisez les pertes de plusieurs créanciers liquidant tour à tour le même collatéral. Que change l'ordre, et que ne change-t-il pas~?
\end{interviewq}
